\documentclass[11pt]{article}
\usepackage{mathblatt}
\usepackage{multicol}


\begin{document}
\pfheftstil
\pfheftkopf{Pythagoras: Seite berechnen \textperiodcentered{} Lösungen}{}
{\small\noindent\textbf{Typische Fehler}\par\smallskip\noindent \textbullet\ Quadrate addiert, obwohl eine Kathete gesucht ist.\par\noindent \textbullet\ falsche Seite als Hypotenuse (längere Kathete; x² = y² + z²).\par\noindent \textbullet\ Wurzel vergessen, c² als Ergebnis.\par\noindent \textbullet\ c = a + b; Wurzel aus jedem Summanden einzeln.\par\noindent \textbullet\ in Figuren: ganze statt halbe Seite, Durchmesser statt Radius, Schräge als Höhe.\par\noindent \textbullet\ Buchstabenfixierung „c ist immer die Hypotenuse“.\par}\medskip
\begin{pfloesung}{}
\lz{1.}{100; 10; 144; 12}{}{}
\end{pfloesung}
\begin{pfloesung}{}
\lz{2.}{ja; ja; nein}{}{}
\end{pfloesung}
\begin{pfloesung}{}
\lz{3.}{Hypotenuse: $41$ cm}{}{}
\end{pfloesung}
\begin{pfloesung}{}
\lz{4.}{rechter Winkel bei C}{ja}{}
\end{pfloesung}
\begin{pfloesung}{}
\lz{5.}{$c = 39$ cm}{Gleichung: $c^2 = 36^2 + 15^2$; Zwischenergebnis: $c^2 = 1\,296 + 225 = 1\,521$; Wurzel: $c = \sqrt{1\,521} = 39$ cm}{}
\end{pfloesung}
\begin{pfloesung}{}
\lz{6.}{$53$ cm}{$28^2 + 45^2 = 2\,809$, Seite }{}
\end{pfloesung}
\begin{pfloesung}{}
\lz{7.}{$\sqrt{97} \approx 9{,}8$ cm}{$c^2 = 97$, $c$}{}
\end{pfloesung}
\begin{pfloesung}{}
\lz{8.}{$130$ cm}{$a = 120$ cm; $c^2 = 16\,900$, $c$}{}
\end{pfloesung}
\begin{pfloesung}{}
\lz{9.}{minus}{Kathete gesucht}{}
\end{pfloesung}
\begin{pfloesung}{}
\lz{10.}{$a = 84$ cm}{Gleichung umstellen: $a^2 = 85^2 - 13^2$; Zwischenergebnis: $a^2 = 7\,225 - 169 = 7\,056$; Wurzel: $a = \sqrt{7\,056} = 84$ cm}{}
\end{pfloesung}
\begin{pfloesung}{}
\lz{11.}{√(10² $-$ 8²) = √36 = 6 cm}{}{}
\end{pfloesung}
\begin{pfloesung}{}
\lz{12.}{Kreuz bei der zweiten Aussage: Das Quadrat über der Hypotenuse ist so groß wie die zwei Quadrate über den Katheten zusammen}{}{}
\end{pfloesung}
\begin{pfloesung}{}
\lz{13.}{h\_s = √(0,68² + 0,40²) = √0,6224 $\approx$ 0,79 m}{}{}
\end{pfloesung}
\begin{pfloesung}{}
\lz{14.}{√(3² + 6²) = √45 = 3√5 cm $\approx$ 6,71 cm}{}{}
\end{pfloesung}
\begin{pfloesung}{}
\lz{15.}{AB = 10√13 $\approx$ 36,1 m}{AB² = 12² + 34² = 1300}{}
\end{pfloesung}
\begin{pfloesung}{}
\lz{16.}{AD = √553 719 $\approx$ 744 m}{AD² = 920² $-$ 541²}{}
\end{pfloesung}
\begin{pfloesung}{}
\lz{17.}{BC = √65 $\approx$ 8,06 m}{BC² = 9² $-$ 4² = 65}{}
\end{pfloesung}
\begin{pfloesung}{}
\lz{18.}{FA = 3√9159 $\approx$ 287,1 m}{FA² = 384² $-$ 255² = 82 431}{}
\end{pfloesung}
\begin{pfloesung}{}
\lz{19.}{h = 3√95 $\approx$ 29,2 cm}{h² = 32² $-$ 13² = 855}{}
\end{pfloesung}
\begin{pfloesung}{}
\lz{20.}{AD = √(14,1² $-$ 7,4²) $\approx$ 12,0 m}{198,81 $-$ 54,76 = 144,05}{}
\end{pfloesung}
\begin{pfloesung}{}
\lz{21.}{√89,96 + 2 $\approx$ 11,5 cm}{Diagonale √(7² + 6,4²) = √89,96 $\approx$ 9,5 cm; Durchmesser 6,4 cm}{}
\end{pfloesung}
\begin{pfloesung}{}
\lz{22.}{Die Summe der Flächeninhalte der Kathetenquadrate ist gleich dem Flächeninhalt des Hypotenusenquadrates}{a² + b² = c²}{}
\end{pfloesung}
\begin{pfloesung}{}
\lz{23.}{z² = x² + y²}{}{}
\end{pfloesung}
\begin{pfloesung}{}
\lz{24.}{z = √(x² + y²)}{}{}
\end{pfloesung}
\begin{pfloesung}{}
\lz{25.}{z = √(x² $-$ y²) (vierte Option)}{z² = x² $-$ y²}{}
\end{pfloesung}
\begin{pfloesung}{}
\lz{26.}{w = √(u² + v²) (dritte Option)}{w² = u² + v²}{}
\end{pfloesung}
\begin{pfloesung}{}
\lz{27.}{Katheten: die Höhe und der Überstand}{}{}
\end{pfloesung}
\begin{pfloesung}{}
\lz{28.}{Radius 3,2 m und Dachhöhe als Katheten, Mantellinie als Hypotenuse $\Rightarrow$ s = √(3,2² + h²)}{}{}
\end{pfloesung}
\begin{pfloesung}{}
\lz{29.}{h = √(193² $-$ 154²) = √13 533 $\approx$ 116,3 m}{}{}
\end{pfloesung}
\begin{pfloesung}{}
\lz{30.}{OR = √(5² + 2²) = √29 $\approx$ 5,39 LE}{}{}
\end{pfloesung}
\begin{pfloesung}{}
\lz{31.}{s = √(1,73² + 3²) $\approx$ 3,46 cm}{}{}
\end{pfloesung}
\begin{pfloesung}{}
\lz{32.}{h\_s = √(3² + 2²) = √13 $\approx$ 3,61 m}{}{}
\end{pfloesung}
\begin{pfloesung}{}
\lz{33.}{s = √(5² + 15,03²) $\approx$ 15,84 cm}{}{}
\end{pfloesung}
\begin{pfloesung}{}
\lz{34.}{√(8,5² + 7,5²) = √128,5 $\approx$ 11,34 m}{}{}
\end{pfloesung}
\begin{pfloesung}{}
\lz{35.}{x = √(38² + 22²) m $\approx$ 43,91 m}{}{}
\end{pfloesung}
\begin{pfloesung}{}
\lz{36.}{25 + √51,87 $\approx$ 32,2 m}{Kegelhöhe = √(8,6² $-$ 4,7²) = √51,87 $\approx$ 7,20 m}{}
\end{pfloesung}
\begin{pfloesung}{}
\lz{37.}{1 Symmetrieachse (BD); BF = 80 $-$ 25 = 55 cm; AB = 5√146 $\approx$ 60,4 cm}{AF = 25 cm; AB² = 25² + 55² = 3650 cm²}{}
\end{pfloesung}
\begin{pfloesung}{}
\lz{38.}{x = 2√7289 $\approx$ 170,8 cm; $\alpha$ = tan$^{-1}$($\tfrac{8}{85}$) $\approx$ 5,4°}{x² = 170² + 16² = 29 156; tan $\alpha$ = 16 : 170 = $\tfrac{8}{85}$ $\approx$ 0,0941}{}
\end{pfloesung}
\begin{pfloesung}{}
\lz{39.}{BC = √20 $\approx$ 4,47; $\beta$ = tan$^{-1}$(2) $\approx$ 63,4°}{AB = 2, AC = 4; BC² = 2² + 4² = 20; tan $\beta$ = 4 : 2 = 2}{}
\end{pfloesung}
\end{document}